\subsection{R-24: exact maximal dimension at the first nondominant size}

\paragraph{Statement R-24.}
The first matrix size beyond the dominance range of R-13 attains
the row--column torus dimension:
\begin{equation}\label{R24:exact}
 \dim\mathcal P_{17,4}=1124=4\cdot17^2-2\cdot17+2.
\end{equation}
Equivalently, the general-section homological defect in R-23 satisfies
$\epsilon_{17}=0$. This is a finite exact result, not an assertion
that $\epsilon_m=0$ for every $m\geq17$.

In particular, the padded pairs $(m,n)=(17,n)$ with
$18\leq n\leq23$ satisfy the strict four-variable dimension
comparison. The pair $(17,23)$ has nonrectangular multiplicity gaps,
including the stretched direction
\[
 \lambda=(41r,17r,17r,17r),\qquad d=4r,
\]
with the effective existence and eventual-stretch conclusions of
G-17/G-18. G-12 excludes rectangular gaps for this pair.

\begin{proof}
The upper bound $1124$ is R-13. We specify a reproducible integer
pencil and a nonzero differential minor for the opposite inequality.
All indices below are zero-based.

Let $s_0=20278274$, and define unsigned $32$-bit integers
\[
 s_{a+1}=1664525s_a+1013904223\pmod{2^{32}},
 \qquad c_a=\left\lfloor s_{a+1}/2^{16}\right\rfloor\bmod101.
\]
For $0\leq i,j<17$, put
\[
 M_{ij}(x_0,x_1,x_2,x_3)
       =\sum_{k=0}^3c_{4(17i+j)+k}x_k.
\]
This recurrence defines all $1156$ integer pencil coefficients
without an external choice.

List the $1140=\binom{20}{3}$ nonnegative integer quadruples
$\alpha=(\alpha_0,\ldots,\alpha_3)$ with sum $17$ in decreasing
lexicographic order. Let $J$ be the following $1140\times1156$
evaluation matrix over $\mathbb F_{101}$:
\begin{equation}\label{R24:evaluation_jacobian}
 J_{\alpha,\,4(17i+j)+k}
 =\alpha_k\,\operatorname{per}
             \big(M(\alpha)\text{ with row }i
                               \text{ and column }j\text{ deleted}\big).
\end{equation}
Retain its first $1124$ rows and all columns except
\begin{equation}\label{R24:omitted_columns}
 \{68i+67:1\leq i\leq15\}
 \ \cup\
 \{1091+4j:0\leq j\leq16\}.
\end{equation}
In increasing row and column order, the resulting square matrix
$J_*$ has
\begin{equation}\label{R24:residue}
 \det J_*=88\pmod{101}.
\end{equation}

Here is the exact computation underlying the finite certificate.
For each numerical matrix $A=M(\alpha)$, the subset recurrence is
\[
 F(\varnothing)=1,\qquad
 F(S)=\sum_{j\in S}A_{|S|-1,j}F(S\setminus\{j\}).
\]
Its reverse sweep starts with $B([17])=1$ and other reverse
coefficients zero. In decreasing subset order, for every $j\in S$
it adds
\[
 B(S)A_{|S|-1,j}\ \text{to }B(S\setminus\{j\}),\qquad
 B(S)F(S\setminus\{j\})\ \text{to }G_{|S|-1,j}.
\]
The resulting $G_{ij}$ are the complementary permanents used in
\eqref{R24:evaluation_jacobian}. All operations are in
$\mathbb F_{101}$. Ordinary exact Gaussian elimination on the
specified square matrix gives \eqref{R24:residue}.

The source file \texttt{check\_permanent\_maximal.cpp} defines the
complete calculation, importing the subset evaluator from the
previous R-13 certificate. The accompanying manifest records the
entire pencil, all evaluation points, selected indices, and raw
matrix. Two additional exact checks were executed: Ryser's
Gray-code formula and its reverse derivative sweep independently
recomputed every one of the $1{,}317{,}840$ matrix entries; a
separate NumPy elimination using unnormalized pivot rows recomputed
the selected determinant. Both agree with \eqref{R24:residue}.
An independent agent then reconstructed the generator and every
evaluation point, checked the compact selected-column description,
and recomputed the increasing-order minor as $88$ by modular
fraction-free Bareiss elimination. That check also verified all
$32$ row--column scaling relations on the full matrix and their
independence. Its scope is the selected minor and the scaling
kernel; the preceding all-entry Ryser calculation is a separate
author check.
These are finite arithmetic checks supporting the stated minor;
they are not an extrapolation to larger matrix sizes.

The matrix $J$ is obtained by evaluating the coefficient
differential of the permanent projection map at the listed points.
Evaluation cannot increase rank. A minor nonzero modulo $101$
is a nonzero integer minor, so the differential over $\mathbb C$
has rank at least $1124$ at this integer pencil. Thus
$\dim\mathcal P_{17,4}\geq1124$, proving
\eqref{R24:exact}; R-23 gives $\epsilon_{17}=0$.

Finally the padded projection image has dimension $1124+3=1127$,
whereas the determinant image for $n\leq23$ has dimension at most
$2\cdot23^2+2=1060$. R-14 and G-17/G-18 apply. At $n=23$,
the padding exponent is $a=6$; the direction
$\nu=(41,17,17,17)$ has $|\nu|=92=23\cdot4$,
$\nu_4>0$, and $\nu_1>4a$. This verifies the strict G-17
conditions for the displayed family. Since $23>4\cdot17/3$,
G-12 excludes rectangles.
\end{proof}

